\chapter{Non-Interactive Compilation}
\label{ch:fs}

The protocols of the previous chapters are public-coin interactive oracle proofs. Their
non-interactive versions replace the verifier by a hash function and the oracles by
Merkle commitments. This chapter proves state-restoration bounds for every binding notion
of \cref{ch:corollaries}, from which the security of the compiled argument in the random
oracle model follows by the transformation of Ben-Sasson, Chiesa and Spooner~\cite{BCS16}.
The proofs use one structural property of the analysis of \cref{ch:path,ch:specialization}:
every bad event is attached to a single challenge and depends only on data fixed before
that challenge, while the evaluation point $c$ of the coupling is never seen by the
prover.

%=====================================================================
\section{The state-restoration game}
%=====================================================================

\begin{definition}[State-restoration game]\label{def:sr}
Fix one of the protocols \cref{prot:sfri,prot:anchored} and one of the winning conditions
of \cref{ch:corollaries} (unique binding, list binding at threshold $\alpha$, or anchored
binding at threshold $\alpha$). A \emph{state-restoration prover} $\mathsf P^\ast$ has
access to a \emph{challenge oracle} that maps every \emph{prefix}, i.e.\ every partial
transcript ending with a prover message, to a challenge for the next round, sampled
uniformly from the appropriate set on the first query and answered consistently
afterwards:
\begin{itemize}
\item a prefix ending with $W_i$ ($0\le i<k$) receives $r_{i+1}\in\K$;
\item for \cref{prot:anchored}, the prefix $W_0$ also receives $z\in\K$;
\item a complete commit transcript $\tau$ receives $s$ independent uniform starting points
$x_0^{(1)},\dots,x_0^{(s)}\in D_0$.
\end{itemize}
$\mathsf P^\ast$ makes at most $Q$ oracle queries and outputs a complete transcript whose
challenges are the oracle answers on its prefixes; the queries needed to verify the output
are counted in $Q$. It wins if the transcript is accepting and the winning condition holds.
The maximal winning probability is denoted $\eps_{\mathrm{SR}}(Q)$.
\end{definition}

The interactive game is the case in which the prover queries each prefix of a single
transcript once, in order.

%=====================================================================
\section{Per-query bad events}
%=====================================================================

Fix $\alpha\in[\rho,1]$, $\delta:=1-\alpha$, and an analysis point $c\in\K^\times$, uniform
and independent of the oracle and of $\mathsf P^\ast$. For a prefix $\tau'$ let
$\Lambda(\tau')$ be the list $\Lambda_\alpha(W_0)$ of the initial word of $\tau'$.

\begin{definition}[Bad answers]\label{def:bad-answer}
An oracle answer is \emph{bad} for its prefix in the following cases.
\begin{enumerate}[label=\textup{(\roman*)}]
\item \emph{MCA.} The prefix ends with $W_i$ and the answer $r_{i+1}$ satisfies
$\varphi_{i+1}(c)\in\Bad\bigl(u[W_i^{(c)}],\delta\bigr)$, where $\varphi_{i+1}(T)$ is the
challenge curve of the protocol ($r_{i+1}T$ or $\zeta_{i+1}+r_{i+1}T$).
\item \emph{Collapse.} The prefix ends with $W_i$ and, for some $G\ne G'$ in
$\Lambda(\tau')$, the polynomial $h:=f_G-f_{G'}\ne0$ satisfies:
$h(\varphi_1(1),\dots,\varphi_i(1),Y_{i+1},\dots,Y_k)\ne0$ but
$h(\varphi_1(1),\dots,\varphi_{i+1}(1),Y_{i+2},\dots,Y_k)=0$, as polynomials in the remaining
variables. \textup{(}For \cref{prot:sfri}, $\varphi_j(1)=r_j$.\textup{)}
\item \emph{Out of domain.} The prefix is $W_0$, the answer is $z$, and two distinct
elements of $\Lambda(W_0)$ take the same value at $z$.
\end{enumerate}
\end{definition}

\begin{lemma}[Per-query probability]\label{lem:per-query}
Each oracle query of a commit round or of the out-of-domain round returns a bad answer with
probability at most
\[
\eps_{\mathrm{cr}}:=\max_{i<k}\eps_{\MCA}(\RS_{i+1},a_{i+1},\delta)
+\binom{\ell_\alpha}2\frac{\max_ia_i-1}{|\K|}+\eps_{\mathrm{OOD}}(\ell_\alpha),
\]
where the last term is present only for \cref{prot:anchored}.
\end{lemma}

\begin{proof}
The prefix, hence $W_i^{(c)}$ and $\Lambda(\tau')$, is fixed when the query is made, and the
answer is uniform and independent of it and of $c$. (i) $\varphi_{i+1}(c)$ is uniform on $\K$
because $c\ne0$, so the probability is at most $\eps_{\MCA}(\RS_{i+1},a_{i+1},\delta)$.
(ii) For each pair, the polynomial before substitution is nonzero and has degree at most
$a_{i+1}-1$ in $Y_{i+1}$; it becomes zero only if $\varphi_{i+1}(1)$ is a common root of its
nonzero coefficients with respect to $Y_{i+1}$, which has probability at most
$(a_{i+1}-1)/|\K|$ since $\varphi_{i+1}(1)$ is uniform. Union over pairs. (iii) is
\cref{lem:ood}.
\end{proof}

\begin{definition}[Exceptional evaluation points]\label{def:zprime}
For a complete commit transcript $\tau$ with challenge curves $\varphi$, let
$Z'(\tau):=\{c'\in\K^\times:\exists G\in\Lambda_\alpha(W_0),\ \gt(c')=f_G(\varphi(c'))\}$ if
$\gt\notin\{f_G(\varphi(T)):G\in\Lambda_\alpha(W_0)\}$, and $Z'(\tau):=\emptyset$ otherwise.
Then $|Z'(\tau)|\le\ell_\alpha\,d$, as in the proof of \cref{thm:symbolic-list}.
\end{definition}

\begin{lemma}[No bad answer implies a small passing ratio]\label{lem:no-bad}
Let $\tau$ be a complete commit transcript, none of whose challenges is a bad answer for its
prefix, such that $c\notin Z'(\tau)$ and the winning condition holds. Then
$\pi^{\mathrm{sym}}_0(\tau)\le\alpha^-$.
\end{lemma}

\begin{proof}
Absence of MCA-bad answers is the event $\neg\mathsf{Bad}^{(c)}$ for $\tau$.
\emph{Case 1:} $\gt\notin\{f_G(\varphi(T)):G\in\Lambda_\alpha(W_0)\}$. If
$\pi^{\mathrm{sym}}_0\ge\alpha$, the argument of \cref{thm:symbolic-list} (via
\cref{lem:spec-passing,lem:master}) gives $c\in Z'(\tau)$, a contradiction; so
$\pi^{\mathrm{sym}}_0<\alpha$, i.e.\ $\le\alpha^-$.
\emph{Case 2:} $\gt=f_G(\varphi(T))$ for some $G\in\Lambda_\alpha(W_0)$. Absence of collapse
answers gives $h(\varphi_1(1),\dots,\varphi_k(1))\ne0$ for every pair, so
$G\mapsto f_G(\varphi(T))$ is injective on the list (the argument of
\cref{prop:curves}(v)) and $G$ is unique. Then no list-binding or unique-binding
condition holds (for unique binding the list is contained in $\{F_\star\}$, so $G=F_\star$).
For anchored binding, $G(z)=f_G(\zeta)=\gt(0)=y$, and absence of out-of-domain-bad answers
makes $G$ the only list element with value $y$ at $z$, i.e.\ $G=G_{z,y}$; so the
anchored condition does not hold either. Hence Case 2 contradicts the winning condition.
\end{proof}

%=====================================================================
\section{The state-restoration bound}
%=====================================================================

\begin{theorem}[State-restoration binding]\label{thm:sr-binding}
For each winning condition of \cref{def:sr} at threshold $\alpha$,
\[
\eps_{\mathrm{SR}}(Q)\le Q\bigl(\eps_{\mathrm{cr}}+(\alpha^-)^s\bigr)+\frac{\ell_\alpha\,d}{|\K|-1}.
\]
For unique binding ($\alpha$ the smallest grid point above $\frac{1+\rho}2$),
\[
\eps_{\mathrm{SR}}(Q)\le Q\Bigl(\frac{(a_1-1)n_1}{|\K|}+\Bigl(\frac{1+\rho}2\Bigr)^s\Bigr)+\frac{d}{|\K|-1}.
\]
\end{theorem}

\begin{proof}
Let $\tau_{\mathrm{out}}$ be the commit transcript of the output and suppose the prover
wins. One of the following holds.
\begin{enumerate}[label=\textup{(\alph*)}]
\item Some challenge of $\tau_{\mathrm{out}}$ is a bad answer for its prefix. Every such
challenge is the answer to one of the $Q$ queries, so by \cref{lem:per-query} this has
probability at most $Q\,\eps_{\mathrm{cr}}$.
\item $c\in Z'(\tau_{\mathrm{out}})$. The whole execution of $\mathsf P^\ast$ and the oracle
is independent of $c$, and $|Z'(\tau_{\mathrm{out}})|\le\ell_\alpha d$; the probability is
at most $\ell_\alpha d/(|\K|-1)$.
\item Otherwise $\pi^{\mathrm{sym}}_0(\tau_{\mathrm{out}})\le\alpha^-$ by \cref{lem:no-bad},
and the starting points returned on $\tau_{\mathrm{out}}$ all lie in
$P_0(\tau_{\mathrm{out}})$. For each queried commit transcript $\tau$ with
$\pi^{\mathrm{sym}}_0(\tau)\le\alpha^-$ this happens with probability at most
$(\alpha^-)^s$, since the answer is fresh; a union bound over the at most $Q$ such queries
gives $Q(\alpha^-)^s$.
\end{enumerate}
For unique binding, $\ell_\alpha\le1$, the collapse and out-of-domain terms vanish, and
$\max_i\eps_{\MCA}(\RS_{i+1},a_{i+1},\delta)\le\max_i(a_{i+1}-1)n_{i+1}/|\K|=(a_1-1)n_1/|\K|$ by
\cref{prop:mca-ud}, since $(a_{i+1}-1)n_{i+1}=n_i-n_{i+1}$ is decreasing in $i$.
\end{proof}

\begin{theorem}[Compiled argument {\cite{BCS16}}]\label{thm:bcs}
Let the oracles be committed with Merkle trees over a random oracle with output length
$\lambda_H$, and derive all challenges by hashing the prefix. For a prover making at most
$Q$ random-oracle queries, the probability of winning any of the binding games of
\cref{ch:corollaries} against the compiled argument is at most
$\eps_{\mathrm{SR}}(Q)+O(Q^2/2^{\lambda_H})$.
\end{theorem}

The transformation theorem of~\cite{BCS16} is stated for soundness; its proof extracts
the committed oracles from the Merkle roots and reduces the compiled adversary to a
state-restoration prover, and applies unchanged to any winning condition that is a
function of the oracle-level transcript.

%=====================================================================
\section{Parameters}
%=====================================================================

% Generated by tools/params (sfri-params). Do not edit by hand.
\begin{table}[ht]\centering\small
\renewcommand{\arraystretch}{1.15}
\begin{tabular}{r|rr|rr|rr}
\hline
 & \multicolumn{2}{c|}{$\F_{p^3}$, $Q=2^{40}$} & \multicolumn{2}{c|}{$\F_{p^3}$, $Q=2^{64}$} & \multicolumn{2}{c}{$\F_{p^4}$, $Q=2^{64}$}\\
$1/\rho$ & $\lambda_{\mathrm{cr}}$ & $s$ & $\lambda_{\mathrm{cr}}$ & $s$ & $\lambda_{\mathrm{cr}}$ & $s$\\
\hline
2 & 132.0 & 406 & 108.0 & --- & 172.0 & 463\\
4 & 131.0 & 249 & 107.0 & --- & 171.0 & 284\\
8 & 130.0 & 203 & 106.0 & --- & 170.0 & 232\\
16 & 129.0 & 186 & 105.0 & --- & 169.0 & 211\\
\hline\end{tabular}
\caption{State-restoration unique binding at $\lambda=128$, $L=2^{20}$, arity $2$. $\lambda_{\mathrm{cr}}=-\log_2(Q\,n_1/|\K|)$ is the security of the commit rounds; $s$ is the smallest query count meeting $2^{-128}$ by \cref{thm:sr-binding}, and --- means that the commit rounds alone exceed $2^{-128}$.}
\label{tab:fs}
\end{table}


\begin{remark}[Reading \cref{tab:fs}]\label{rem:fs-reading}
\begin{enumerate}[label=\textup{(\roman*)}]
\item The commit rounds contribute $Q\,n_1/|\K|$, which does not decrease with $s$. Over
$\K=\F_{p^3}$ with $Q=2^{64}$ this is about $2^{-107}$ at $\rho=\frac14$, so $128$-bit
non-interactive unique binding is unreachable for every $s$; it is reachable with
$Q\le2^{40}$, or with $\K=\F_{p^4}$.
\item Where reachable, the query count is
$s=\lceil(128+\log_2Q)/\log_2\frac2{1+\rho}\rceil$ up to the small correction of the other
terms; for $\F_{p^4}$ and $Q=2^{64}$ this gives $463,284,232,211$.
\item Proof-of-work on the query-round prefix with $w$ bits multiplies the query-round term
by $2^{-w}$ and leaves the commit term unchanged.
\item The anchored and list-binding variants have the same structure with $\alpha^-$ in
place of $\frac{1+\rho}2$; their commit term contains $\eps^J_{\MCA}$ and is subject to the
same field-size condition.
\end{enumerate}
\end{remark}
